% 06d-engineering.tex 中文对照版（译自 sections/06d-engineering.tex；不加 \label）
\subsection{工程验证与系统可靠性}

契约的可信度只能等于执行它的测试的可信度，而断言数量本身不是证据。因此我们从三条与计数无关的轴线审计测试套件：变异测试（它能发现被注入的语义故障吗？）、断言强度分级（断言有区分力吗，还是同义反复或自指？），以及实测覆盖率。

\paragraph{变异测试：契约是活的吗？} 如果测试永远不会失败，契约就只是一份文档，而不是可执行工件；因此审计注入了 44 个单故障语义变异体——没有语法错误或导入错误，全部能加载并执行——分布于 JavaScript 插件（29 个）、Python 引擎（11 个）与规则文件（4 个）。每个变异体单独施加，并运行\emph{完整}测试套件：每轮 1{,}249 条断言，约 55{,}000 次断言执行。只有断言失败才计为杀死，从不因进程崩溃而计；没有任何变异体是被错误杀死的。审计时套件杀死 $38/44=86.4\%$（JavaScript $24/29=82.8\%$，Python $10/11=90.9\%$，规则 $4/4=100\%$）。在 40 个代码变异体中，34 个被行为断言杀死（$85.0\%$）；4 个规则变异体是由跨文件一致性断言而非行为断言捕获的，我们保留这一区分，使其可见。下面四个盲点被最小回归断言钉住之后，套件杀死 $42/44=95.5\%$，并且杀死 $42/42=100\%$ 的\emph{非等价}变异体。分数能从 86.4\% 升到 95.5\%——这正是“契约是活的、而不只是存在”的具体含义。

\paragraph{对抗式流程。} 该审计在设计上就是对抗式的。变异测试由一名审查者执行，断言与覆盖率分析由第二名审查者执行，另有第三名审查者在另一份副本上独立复跑其中一个子集。全部工作都在副本上进行；30 个实现文件在审计前后做了哈希，未发生变化（校验 $=30$，变动 $=0$）。只有测试文件被改动。

\paragraph{四个真实的盲点。} 存活下来的变异体暴露了四个真实盲点，且都位于同一结构位置：单个写批次内部的一致性。两个与取代有关——删除工作副本失效逻辑后，当一个批次内两个新值替换同一旧条目时会留下两个墓碑（tombstone）；而颠倒墓碑顺序，则会让已被取代的值在链式批次中依 last-write-wins 回到 \emph{active}。一个与冲突检测有关：其 active 状态过滤器是端到端测试永远到不了的防御层，因为检索阶段已经滤过，只有直接的单元测试才能钉住它。一个与质量标签闭集有关：每一条标签测试句都同时匹配第二个关键词，于是删掉一个关键词后套件仍然全绿；修复办法是把每个标签钉到单一关键词，并断言被引用的证据是逐字子串。以上每项都被复现、用一条最小回归断言钉住（已验证它在原始代码上通过、在变异体上失败），并加入仓库（22 条 Python、15 条 JavaScript）。剩余两个存活者是等价变异体：积分钳制的负分支在不变量 $\mathit{integral}\ge 0$ 下不可达，而该不变量对契约的正增益成立，因此 $42/42=100\%$ 的\emph{非等价}变异体被杀死。我们只报告这一句，而不是笼统声称“完全检出”。

\paragraph{断言强度。} 对 238 条断言做固定种子抽样（165 条 Python、73 条 JavaScript）并启发式分级：179 条强（$75.2\%$）、59 条弱（$24.8\%$），样本中零同义反复。对 1{,}192 条静态断言做完整启发式扫描，Python 评为 597 强、227 弱、1 条同义反复；JavaScript 为 278 强、89 弱，即总体上 $26.6\%$ 为弱或同义反复。我们发现并修复了一条同义反复断言（导入语句已由模块头执行过）、两条 oracle 自指断言（其期望值是从实现自身的 unsafe-control-action 表里读的；已改为对照测试侧的字面量表重写，该文件由 159 条断言增至 161 条），以及一条阈值输入取自被审计配置的断言。每次重写都通过注入对应故障并观察新断言失败来验证。因此，我们不声称 892 条断言就是 892 条独立、精确的字面量；我们声称：样本中没有同义反复；完整扫描发现了一条（已修）；大约四分之一的断言被评为弱。

\paragraph{异常路径。} 在严格口径下（只统计非法、空、超大、注入、超时、格式错误或边界输入），异常路径占比为 Python $199/825=24.1\%$、JavaScript $82/367=22.3\%$（放宽口径的上界分别为 $43.3\%$ 与 $54.2\%$）。我们为自己设定的“要求补健壮性测试”的门槛是 $20\%$；两个套件都越过了，但只是刚刚越过；若剔除阈值边界断言，Python 会降到约 $11\%$。准确的说法是“刚过线，并不宽裕”。

\paragraph{覆盖率。} 覆盖率是实测而非估算：Python 使用 PEP~669 \texttt{sys.monitoring}，JavaScript 使用 V8 覆盖数据。对 \texttt{engine/}，行覆盖为 $89.95\%$、分支覆盖为 $80.71\%$；\texttt{control.py} 达到 $141/141$ 行与 $56/56$ 分支。全仓库 Python 覆盖为行 $45.81\%$、分支 $43.17\%$，被 \texttt{tools/} 拉低（$32.92\%$，26 个文件中只有 3 个被执行）。我们不用覆盖率论证死代码不存在：在被测配置下 \texttt{control.py} 没有被执行的行，却仍有两个等价变异体存活——这正是覆盖率单独无法确立测试强度的原因。

\paragraph{基线记录。} 审计基线为 13 个 Python 文件、868 条断言，以及 381 条 JavaScript 断言；审计新增 22 条与 2 条 Python 断言、15 条 JavaScript 断言，最终为 14 个文件、892 与 396 条断言，全部通过，实现源码未改动。

\paragraph{局限。} 四条局限限定了上述主张。仓库没有版本控制历史，因此我们能描述当前状态，但无法证明任何断言从未被弱化。变异为单故障且未覆盖每个模块。不存在跨语言差分测试：没有任何变异体同时被两个套件杀死，因此单侧的语义漂移只对它自己的套件可见。最后，\texttt{tools/} 的覆盖率是下界：以子进程方式调用的工具会被记录为未执行。
